% Figure 1 Caption: LMSYS Holdout Semantic Structure

\begin{figure}[t]
\centering
\includegraphics[width=0.95\textwidth]{experiments_v1/01_figure/results/figure1_lmsys_holdout_pca.png}
\caption{\textbf{Proof of Learnability: LLM routing exhibits bimodal semantic structure.} 
We project 1,871 LMSYS prompts (dev + holdout sets) onto the first two principal components of a 
32-component PCA model (PC1: 3.10\%, PC2: 2.29\%, cumulative: 5.39\% variance). \textbf{Left:} 
Semantic scatter plot colored by PC1 position reveals two spatially distinct clusters. The Low PC1 
cluster (blue, 82.4\%, 1,541 prompts) represents routine semantic tasks where mid-tier models perform 
adequately. The High PC1 cluster (red, 17.6\%, 330 prompts) isolates complex coding and reasoning 
tasks requiring flagship capabilities. The dashed line at PC1 $= 0.3$ marks the cluster separation 
boundary. Kernel density estimation (KDE) contours (blue) highlight the spatial coherence of the 
Low PC1 cluster. \textbf{Right:} The spatial cluster distribution demonstrates the economic 
opportunity: 82.4\% of production traffic occupies regions where expensive flagship models are 
unnecessary. This asymmetric distribution provides the empirical foundation for cost-effective 
contextual routing: by exploiting the PC1 decision boundary, our approach targets 80\%+ cost 
reduction on routine tasks while preserving quality on the critical 17.6\% reasoning cluster.}
\label{fig:lmsys_holdout_structure}
\end{figure}

